\section*{Step 261: RH Attack Foreclosure Conjecture}

\paragraph{Candidate Meta-Theorem.}
Under the Six Birds Foundations III typed-condition discipline, no
Riemann-Hypothesis proof can be constructed by carrier pivoting, bridge import,
or cascade-internal computation alone, in isolation from named external
classical content.

\paragraph{Foreclosed Internal Moves.}
\begin{enumerate}
\item \textbf{Carrier pivoting.}  Choosing a different RH-equivalent carrier is
governed by the Carrier Dichotomy and CRCFT modes.  A CRE carrier is classified
as target-equivalent, carrier-typed terminal, or bridge-failed.  None supplies
non-circular internal closure.
\item \textbf{Bridge import.}  Composing a proved adjacent theorem with a bridge
to RH is governed by Bridge Impossibility.  The composite becomes
RH-equivalent, so the bridge itself is RH-strength.
\item \textbf{Cascade-internal computation.}  Extending reductions is governed
by CTMT and CTMT recursion.  Resolving one stuck gate typically exposes a deeper
gate or a precision/theorem-unstated terminus.
\end{enumerate}

\paragraph{Coverage.}
The nine current candidate findings cover the standard internal move families:
carrier pivots, bridge imports, cascade computations, per-$L$ automorphic RH
analogues, coefficient and zero correlations, subconvexity and moment growth,
and zero-density/sieve-type estimates.

\paragraph{Not Foreclosed.}
The conjecture does not foreclose a proof of RH.  It preserves four genuine
interfaces: resolution of named external classical theorems, a fresh typed
mechanism beyond the current findings, a fresh carrier outside the surveyed
atlas, and direct verification using external closed-form content.

\paragraph{Status.}
Candidate framework meta-theorem; corpus-pending; awaiting cross-track
replication as a methodology test.
